% Image credits for the photographs and AI illustrations of Book 1.
% Machine-readable records live in images/book1/CREDITS.md; keep the two in
% sync. The Book 1 writer owns this file: add one \item per photograph (by
% chapter) as photographs land, and remove the "(none yet)" item with the first.
% If the book ends with no photograph, the writer keeps the first paragraph only.
\chapter*{Image Credits}

The drawings, diagrams and charts in this book are original. Each chart is
redrawn from published data by a tested script and names its source under it;
the datasets are credited in full in the Data Credits.
Photographs are used under their respective free licenses, with thanks to
their authors; full source links and license URLs for every photograph are
listed in \texttt{images/book1/CREDITS.md} in the book's repository.

\bigskip
Some figures are AI-generated illustrations of everyday scenes, produced with
OpenAI image generation from prompts written by the book's editors, and
reviewed for accuracy before inclusion. They show generic scenes only, never a
real person, a real product or a real brand. The full prompt for every
generated image is recorded in \texttt{images/book1/ai/PROMPTS.md} in the
repository.

\bigskip
\noindent Photographs, by chapter:
\begin{itemize}\small
  \item Chapter~\ref{ch:g6:school-and-work}: boys who sorted coal at a mine in
    Pennsylvania, 1911. Photograph by Lewis Hine for the National Child Labor
    Committee; Library of Congress, Prints and Photographs Division,
    \texttt{nclc.01138}; public domain.
  \item Chapter~\ref{ch:g6:a-connected-world}: aerial view of the container port
    of Miami. Photograph by James R. Tourtellotte, U.S. Customs and Border
    Protection; public domain (Wikimedia Commons).
  \item Chapter~\ref{ch:g7:specialisation}: a pin-making workshop, plate II of
    ``\'Epinglier'' in the plates of the \emph{Encyclop\'edie} (1762), drawn by
    Goussier and engraved by Defehrt; public domain (Wikimedia Commons).
  \item Chapter~\ref{ch:g7:money}: an electrum coin (trite) of Sardis, Lydia.
    Biblioth\`eque nationale de France, Gallica \texttt{btv1b8502096t}; public
    domain (Wikimedia Commons).
  \item Chapter~\ref{ch:g8:production-and-productivity}: a car assembly line at
    Highland Park, Michigan, 1913. Library of Congress, Prints and Photographs
    Division, \texttt{LCCN 2011661021}; public domain (Wikimedia Commons).
  \item Chapter~\ref{ch:g8:inflation-and-price-indices}: a banknote of 100 thousand
    million marks, Berlin, 26 October 1923. National Numismatic Collection, National
    Museum of American History (image by Godot13); public domain (Wikimedia Commons).
  \item Chapter~\ref{ch:g9:international-trade}: a container ship being loaded in
    the port of Barcelona, 2010. Photograph by Ivan Mlinaric; CC BY 2.0
    (Wikimedia Commons).
  \item Chapter~\ref{ch:g9:institutions-and-prosperity}: the Korean peninsula at
    night, 30 January 2014, photographed from the International Space Station by
    an Expedition 38 crew member. NASA; public domain (Wikimedia Commons).
\end{itemize}
